\section{Related Work}
\label{sec:related_work}

Recent audio deepfake resources have developed along four complementary directions: challenge benchmarks, diversity-oriented datasets, delivery-oriented robustness studies, and reproducibility frameworks. Rather than comparing only dataset scale, Table~\ref{tab:rw_compare} focuses on the aspects most relevant to this work: delivery scope, structured composition, and matched-parent control.

\begin{table*}[t]
\centering
\footnotesize
\setlength{\tabcolsep}{3pt}
\renewcommand{\arraystretch}{1.01}
% \caption{Representative prior resources versus ChainBench-ADD. \cmark denotes explicit support, (\cmark) partial or special-case support, \textemdash\ absence, and \textsc{n/a} not applicable because the work is primarily an evaluation framework.}
\caption{Representative prior resources versus ChainBench-ADD. \cmark\ denotes explicit support, (\cmark) partial support, \textemdash\ absence, and \textsc{n/a} not applicable because the work is primarily an evaluation framework.}
\label{tab:rw_compare}
\vspace{-0.5em}
\begin{tabular}{@{}p{2.95cm}p{2.45cm}p{4.85cm}cc@{}}
\toprule
\textbf{Resource} & \textbf{Focus} & \textbf{Post-generation delivery} & \textbf{Ordered comp.} & \textbf{Matched-parent control} \\
\midrule
ASVspoof 2021~\cite{liu2023asvspoof} & challenge benchmark & codec/channel/room effects in challenge tracks & \textemdash & \textemdash \\
ADD 2023~\cite{yi2024add} & task benchmark & not delivery-centric & \textemdash & \textemdash \\
CFAD~\cite{ma2024cfad} & generator/domain range & noise + codec condition sets & \textemdash & \textemdash \\
CodecFake~\cite{xie2025codecfake} & ALM/codec coverage & generator-side codec effects & \textemdash & \textemdash \\
ADD-C~\cite{shi2025benchmarking} & comms robustness & codec + packet-loss conditions & (\cmark) & \textemdash \\
ReplayDF~\cite{muller2025replay} & replay robustness & playback + re-recording pipeline & (\cmark) & \textemdash \\
AUDDT~\cite{zhu2025auddt} & evaluation toolkit & \textsc{n/a} & \textsc{n/a} & \textsc{n/a} \\
Spoof-SUPERB~\cite{ali2026superb} & SSL benchmark & \textsc{n/a} & \textsc{n/a} & \textsc{n/a} \\
\textbf{ChainBench-ADD (ours)} & \textbf{structured delivery} & \textbf{multi-family post-generation delivery} & \textbf{\cmark} & \textbf{\cmark} \\
\bottomrule
\end{tabular}
\end{table*}

Early evaluation standards were shaped largely by the ASVspoof series. ASVspoof 2021~\cite{liu2023asvspoof} formalized logical-access, physical-access, and deepfake evaluation while incorporating compression, transmission, and acoustic propagation effects; ASVspoof 5~\cite{wang2026asvspoof} further broadened speaker, acoustic, and attack diversity. In parallel, ADD 2023~\cite{yi2024add} expanded evaluation beyond utterance-level bona fide/spoof discrimination to include manipulation localization and algorithm recognition, and PartialSpoof~\cite{zhang202partial} focused on partially manipulated utterances. Earlier resources such as WaveFake~\cite{frank2021wavefake} and FakeAVCeleb~\cite{khalid2021fakeavceleb} also helped establish benchmarked evaluation for synthetic and multimodal deepfakes.

A second line of work emphasizes broader generator, language, and domain coverage. In-the-Wild~\cite{muller2022generalize} highlighted the gap between laboratory benchmarks and real online media. MLAAD~\cite{muller2024mlaad} expanded multilingual and multi-generator coverage, LibriSeVoc~\cite{sun2023ai} emphasized vocoder artifacts and self-vocoding conditions, CFAD~\cite{ma2024cfad} introduced a Chinese benchmark under noisy and transcoding conditions, and CodecFake~\cite{xie2025codecfake} focused on codec-based and neural-codec generation. More recent datasets~\cite{huang-etal-2025-speechfake,mawalim2025jmad,ciobanu2026xmad,wang2025audeter} continue this trend toward larger-scale multilingual, cross-domain, and open-world evaluation. Despite that, the main axis of variation in these resources remains the generator, language, or source domain.

More closely related to our setting are delivery-oriented robustness studies. ADD-C~\cite{shi2025benchmarking} shows that codec compression, packet loss, and communication-channel effects can noticeably affect detector performance in VoIP- and telephony-like scenarios, while ReplayDF~\cite{muller2025replay} reports similar degradation under playback and re-recording. Related discussion in~\cite{delgado2026presentation} further points to a realism gap between raw generated audio and audio after practical presentation through communication channels or platform pipelines. These studies establish delivery as an important source of performance variation, but they typically examine communication, replay, or platform effects separately rather than as parts of one structured delivery process.

Recent frameworks have also improved reproducibility at the system level. AUDDT~\cite{zhu2025auddt} standardizes cross-dataset evaluation across public corpora, and Spoof-SUPERB~\cite{ali2026superb} provides a SUPERB-style comparison of self-supervised speech representations for audio deepfake detection. These efforts improve fair and reproducible comparison, but they primarily reorganize or evaluate existing resources rather than introduce a matched-parent benchmark centered on post-generation delivery.

ChainBench-ADD is designed to complement this landscape by modeling post-generation delivery as an ordered process under matched-parent control. To the best of our knowledge, existing representative resources have not jointly provided multiple post-generation delivery families, explicit ordered templates, and matched-parent local analysis within one unified benchmark.
